\chapter{Annexes}
\label{ch:annexes}

\section{Descriptions complémentaires des cas d'utilisation}

Cette annexe regroupe des éléments complémentaires utiles à la lecture du rapport : correspondance avec
les cas d'utilisation, emplacements réservés aux visuels, commandes de lancement, commandes de test,
résultats de validation et exemples de structure CSV.

\begin{longtable}{L{4cm}L{10cm}}
    \caption{Cas d'utilisation et éléments de validation}\\
    \toprule
    \textbf{Cas d'utilisation} & \textbf{Complément attendu} \\
    \midrule
    \endfirsthead
    \toprule
    \textbf{Cas d'utilisation} & \textbf{Complément attendu} \\
    \midrule
    \endhead
    Authentification & Décrite par un cas d'utilisation textuel et un diagramme de séquence JWT au chapitre Conception. \\
    Recherche de pharmacie & Décrite par un cas d'utilisation textuel, un diagramme de séquence et une vue mobile de recherche. \\
    Consultation des gardes & Décrite dans les besoins fonctionnels, les intégrations API et la vue mobile du calendrier. \\
    Gestion des pharmacies & Décrite dans les besoins administrateur, la réalisation web et la vue de gestion des pharmacies. \\
    Import CSV & Décrit dans la réalisation backend, la stratégie de validation et les exemples de structure CSV. \\
    Audit logs & Décrit dans la sécurité, les tests backend et les indicateurs du tableau de bord. \\
    \bottomrule
\end{longtable}

\section{État des captures et diagrammes}

\begin{longtable}{L{5cm}L{9cm}}
    \caption{Éléments visuels à insérer}\\
    \toprule
    \textbf{Élément} & \textbf{Statut} \\
    \midrule
    \endfirsthead
    \toprule
    \textbf{Élément} & \textbf{Statut} \\
    \midrule
    \endhead
    Diagramme de cas d'utilisation & Emplacement réservé ; insérer le diagramme final exporté. \\
    Diagramme de classes & Emplacement réservé ; insérer le diagramme final exporté. \\
    Diagramme de séquence authentification JWT & Emplacement réservé ; insérer le diagramme final exporté. \\
    Diagramme de séquence recherche pharmacie & Emplacement réservé ; insérer le diagramme final exporté. \\
    Diagramme d'activité import CSV & Emplacement réservé ; insérer le diagramme final exporté. \\
    Diagramme de composants & Emplacement réservé ; insérer le diagramme final exporté. \\
    Diagramme de déploiement & Emplacement réservé ; insérer le diagramme final exporté. \\
    Planning de réalisation & Emplacement réservé ; insérer le planning final exporté. \\
    Capture tableau de bord administrateur & Emplacement réservé ; insérer la capture réelle. \\
    Capture gestion des pharmacies & Emplacement réservé ; insérer la capture réelle. \\
    Capture accueil mobile & Emplacement réservé ; insérer la capture réelle. \\
    Capture calendrier des gardes & Emplacement réservé ; insérer la capture réelle. \\
    Capture catalogue médicaments & Emplacement réservé ; insérer la capture réelle. \\
    \bottomrule
\end{longtable}

\section{Commandes de lancement}

\subsection*{Backend}

\begin{lstlisting}
cd backend_pharmacie
python -m venv venv
source venv/bin/activate
pip install -r requirements.txt
uvicorn main:app --reload
\end{lstlisting}

\subsection*{Interface d'administration}

\begin{lstlisting}
cd admin_pharmacie
npm install
npm run dev
\end{lstlisting}

\subsection*{Application mobile}

\begin{lstlisting}
cd ouerkema-pharmacieconnect-4c94773cce7d
npm install
npm start
\end{lstlisting}

\section{Commandes de test}

\subsection*{Backend}

\begin{lstlisting}
cd backend_pharmacie
source venv/bin/activate
pytest
\end{lstlisting}

\begin{lstlisting}
python -m pytest tests/test_analytics_dashboard.py \
  tests/test_assistant_regions.py \
  tests/test_audit_logs.py --tb=short
\end{lstlisting}

\subsection*{Administration web}

\begin{lstlisting}
cd admin_pharmacie
npm run lint
npm run build
\end{lstlisting}

\subsection*{Application mobile}

\begin{lstlisting}
cd ouerkema-pharmacieconnect-4c94773cce7d
npm test
npm run lint
\end{lstlisting}

\section{Résultats de validation du 13 mai 2026}

\begin{longtable}{L{4cm}L{5cm}L{5cm}}
    \caption{Résultats détaillés des commandes exécutées}\\
    \toprule
    \textbf{Commande} & \textbf{Objectif} & \textbf{Résultat} \\
    \midrule
    \endfirsthead
    \toprule
    \textbf{Commande} & \textbf{Objectif} & \textbf{Résultat} \\
    \midrule
    \endhead
    \texttt{pytest} ciblé backend & Vérifier les indicateurs, l'audit et les restrictions régionales & 14 tests réussis, 21 avertissements, 4,45 s. \\
    \texttt{npm run build} & Vérifier la génération de l'administration web & Build réussi, artefacts générés dans \texttt{dist/}. \\
    \texttt{npm test -- --runInBand} & Vérifier les tests unitaires mobiles & 7 tests réussis sur la recherche mobile par localisation. \\
    \texttt{pytest} upload médicaments & Vérifier l'import des médicaments & Timeout après 20 s ; correction ou isolation nécessaire avant validation exhaustive. \\
    \bottomrule
\end{longtable}

\section{Variables d'environnement}

\begin{table}[H]
    \centering
    \caption{Variables d'environnement utilisées par le backend}
    \begin{tabular}{L{6.2cm}L{6.8cm}}
        \toprule
        \textbf{Variable} & \textbf{Rôle} \\
        \midrule
        \texttt{DATABASE\_URL} & Chaîne de connexion à la base de données. \\
        \texttt{SECRET\_KEY} & Clé de signature des jetons JWT. \\
        \texttt{ALGORITHM} & Algorithme de signature des jetons. \\
        \texttt{ACCESS\_TOKEN\_EXPIRE\_MINUTES} & Durée de validité du jeton d'accès. \\
        \texttt{REFRESH\_TOKEN\_EXPIRE\_DAYS} & Durée de validité du jeton de rafraîchissement. \\
        \texttt{FRONTEND\_URL} & Origine autorisée pour l'interface web. \\
        \texttt{TRUSTED\_HOSTS} & Hôtes autorisés par le backend. \\
        \texttt{REDIS\_URL} & URL du service Redis lorsque le cache est utilisé. \\
        \texttt{ENABLE\_REDIS\_CACHE} & Active ou désactive le cache Redis si configuré. \\
        \bottomrule
    \end{tabular}
\end{table}

\section{Exemples CSV}

\subsection*{Pharmacies}

\begin{table}[H]
    \centering
    \caption{Exemple de structure CSV pour les pharmacies}
    \begin{tabular}{L{4cm}L{9cm}}
        \toprule
        \textbf{Colonne} & \textbf{Description} \\
        \midrule
        \texttt{name} & Nom de la pharmacie. \\
        \texttt{latitude} & Latitude GPS. \\
        \texttt{longitude} & Longitude GPS. \\
        \texttt{address} & Adresse complète ou approximative. \\
        \texttt{phone} & Numéro de téléphone. \\
        \texttt{governorate} & Gouvernorat utilisé pour le filtrage régional. \\
        \bottomrule
    \end{tabular}
\end{table}

\subsection*{Médicaments}

\begin{table}[H]
    \centering
    \caption{Exemple de structure CSV pour les médicaments}
    \begin{tabular}{L{5.2cm}L{7.8cm}}
        \toprule
        \textbf{Colonne} & \textbf{Description} \\
        \midrule
        \texttt{code\_pct} & Identifiant métier du médicament. \\
        \texttt{nom\_commercial} & Nom commercial. \\
        \texttt{prix\_public\_dt} & Prix public en dinars. \\
        \texttt{tarif\_reference\_dt} & Tarif de référence. \\
        \texttt{categorie\_remboursement} & Catégorie de remboursement. \\
        \texttt{dci} & Dénomination commune internationale. \\
        \texttt{ap} & Autorisation ou indicateur AP selon la source de données. \\
        \bottomrule
    \end{tabular}
\end{table}
